\documentclass[11pt,a4paper]{article}

% --- Packages ---------------------------------------------------------------
\usepackage[utf8]{inputenc}
\usepackage[T1]{fontenc}
\usepackage[margin=2.5cm]{geometry}
\usepackage{amsmath,amssymb,amsthm}
\usepackage{mathtools}
\usepackage{booktabs}
\usepackage{array}
\usepackage{tabularx}
\usepackage{longtable}
\usepackage{multirow}
\usepackage{graphicx}
\usepackage{caption}
\usepackage{subcaption}
\usepackage{enumitem}
\usepackage{multicol}
\usepackage{hyperref}
\usepackage{xcolor}
\usepackage{listings}
\usepackage{microtype}
\usepackage{fancyhdr}
\usepackage{authblk}
\usepackage{siunitx}
\usepackage{titlesec}

\hypersetup{
    colorlinks=true,
    linkcolor=black,
    citecolor=blue!60!black,
    urlcolor=blue!60!black,
    pdftitle={Trend Following on BIST100: A Low-Frequency Momentum Strategy},
    pdfauthor={EC581 Project}
}

\sisetup{
    detect-all,
    group-separator={,},
    group-minimum-digits=4
}

\lstset{
    basicstyle=\ttfamily\small,
    breaklines=true,
    backgroundcolor=\color{gray!8},
    frame=single,
    rulecolor=\color{gray!30},
    keywordstyle=\color{blue!60!black}\bfseries,
    commentstyle=\color{green!50!black}\itshape,
    stringstyle=\color{red!60!black},
    showstringspaces=false,
    language=Python
}

\pagestyle{fancy}
\fancyhf{}
\rhead{\small EC581 -- Trend Following Project}
\lhead{\small BIST100 Low-Frequency Momentum}
\rfoot{\thepage}

\titleformat{\section}{\Large\bfseries}{\thesection}{0.6em}{}
\titleformat{\subsection}{\large\bfseries}{\thesubsection}{0.5em}{}

% --- Title ------------------------------------------------------------------
\title{\textbf{A Low-Frequency Momentum Trading Strategy \\
\large on Turkish Equities (BIST100)} \\[0.5em]
\large EC581 -- Algorithmic Trading and Quantitative Strategies \\[0.2em]
\normalsize Project Section 4: Trend Following}

\author[]{Bogazici University -- Course Project}
\affil[]{Instructor: Dr.~Ayhan Y\"uksel}
\date{May 2026}

% ============================================================================
\begin{document}
% ============================================================================

\maketitle
\thispagestyle{empty}

\begin{abstract}
\noindent We design, implement and evaluate a long/short trend-following strategy for
Turkish equities (BIST100) inspired by Harris and Yilmaz (2008), \emph{A Momentum
Trading Strategy Based on the Low Frequency Component of the Exchange Rate}.
Four trend extractors --- EMA crossover (S1), EMA direction (S2), Hodrick--Prescott
filter direction (S3), and LOWESS direction (S4) --- are tuned in-sample on the
BIST100 Index (Phase~1) and then deployed on 29 constituent stocks (Phase~2) with
an optional BIST100-HP regime gate. Statistical significance is established with a
run-length-matched Monte-Carlo bootstrap ($N{=}1000$). The Hodrick--Prescott
direction rule (S3) is retained as the principled primary strategy by virtue of its
theoretical alignment with Harris--Yilmaz; LOWESS (S4) is carried as the
\emph{empirical} benchmark and reported head-to-head.
On the BIST100 Index over 23.3~years, S3~HP reaches Sharpe $0.82$ ($p<0.001$
vs.~matched-random) and S4~LOWESS reaches Sharpe $0.88$ ($p<0.001$). The
equal-weight 29-stock portfolio gives S3~HP a Sharpe of $0.90$ (CAGR $15.3\%$,
MDD $-26.8\%$) and S4~LOWESS a Sharpe of $1.05$ (CAGR $18.2\%$, MDD $-28.5\%$),
compared to BIST100 buy-and-hold Sharpe $0.92$ (CAGR $23.5\%$, MDD $-63.4\%$).
The regime filter consistently acts as a \emph{variance reducer} rather than a
Sharpe enhancer. We close with a detailed methodological self-critique and a
set of concrete improvement directions covering universe construction, signal
combination, position sizing, transaction-cost realism, and statistical testing.
\end{abstract}

\tableofcontents
\newpage

% ============================================================================
\section{Introduction}
% ============================================================================

\subsection{Project Scope}

This report documents the project for Section~4 (\emph{Trend Following}) of the
EC581 course brief at Bogazici University. The anchoring article is Harris and
Yilmaz (2008/2009), \emph{A Momentum Trading Strategy Based on the Low Frequency
Component of the Exchange Rate}, published in the \emph{Journal of Banking and
Finance}. The brief mandates the following constraints, all of which are
respected:

\begin{itemize}[leftmargin=1.2em,itemsep=2pt]
    \item Backtest engine: \textbf{Backtrader}.
    \item Initial capital: \textbf{1{,}000{,}000~TL}.
    \item Per-trade notional: \textbf{100{,}000~TL}, allocated via the
          course-supplied \texttt{FixedCashSizer}.
    \item Commissions and slippage: \textbf{zero}.
    \item Order type (base strategy): \texttt{bt.Order.Market}.
    \item Frequency: \textbf{daily} bars.
    \item Universe Phase~1: \textbf{BIST100 Index} for hyperparameter selection.
    \item Universe Phase~2: \textbf{BIST100 constituents} (current membership;
          survivorship bias acknowledged but uncorrected per brief).
    \item Significance: \textbf{Monte-Carlo $p$-value} on the Sharpe ratio
          against random strategies matching trade frequency, holding period and
          long/short proportion, with $N\!\approx\!1000$.
\end{itemize}

The project is to (i)~implement the four trend extractors described in
\S\ref{sec:strategies}, (ii)~tune each on the index, (iii)~pick a primary,
(iv)~deploy it on individual stocks both in a base and a regime-gated form, and
(v)~test statistical significance throughout.

\subsection{Why Trend Following?}

Classical efficient-market theory predicts no exploitable serial dependence in
returns. In practice, persistent autocorrelation in returns has been documented
across asset classes and decades (Moskowitz, Ooi and Pedersen, 2012; Asness,
Moskowitz and Pedersen, 2013). Behavioural explanations --- anchoring,
under-reaction to news, herding --- and structural ones --- slow capital flows,
leverage cycles, central-bank reaction functions --- all predict that prices
adjust to information gradually rather than instantaneously. A strategy that is
\emph{long when the slow trend is up and short when it is down} aims to harvest
that drift.

\subsection{Why Low-Frequency Filters?}

Harris and Yilmaz~(2008) decompose daily FX series into a low-frequency
component (extracted with the Hodrick--Prescott filter from business-cycle
econometrics) and a high-frequency residual. Their three operationally useful
claims are:

\begin{enumerate}[leftmargin=1.5em,itemsep=2pt]
    \item The low-frequency component is \emph{highly persistent}: once its slope
          flips sign, it tends to keep that sign for many periods.
    \item Rules based on the \textbf{sign of the change in the low-frequency
          component} outperform classical moving-average crossover rules
          out-of-sample across major currency pairs.
    \item Low-pass filters, while non-causal during fitting, produce smoother
          right-edge estimates than moving averages of comparable lag, reducing
          whipsaws.
\end{enumerate}

We test these claims on Turkish equities.

% ============================================================================
\section{Theoretical Background}\label{sec:theory}
% ============================================================================

\subsection{The Hodrick--Prescott Filter}

Given a price series $\{P_t\}_{t=1}^{T}$, the HP decomposition writes
$P_t = \tau_t + c_t$, where $\tau_t$ is a smooth trend and $c_t$ the residual
cycle. The trend solves
\begin{equation}
\widehat{\boldsymbol{\tau}} \;=\; \arg\min_{\boldsymbol{\tau}}\;
\underbrace{\sum_{t=1}^{T}(P_t - \tau_t)^2}_{\text{fit}}
\;+\;\lambda\,
\underbrace{\sum_{t=2}^{T-1}\bigl((\tau_{t+1}-\tau_t)-(\tau_t-\tau_{t-1})\bigr)^2}_{\text{smoothness}},
\label{eq:hp}
\end{equation}
with closed form
\begin{equation}
\widehat{\boldsymbol{\tau}} \;=\; (\mathbf{I} + \lambda\, \mathbf{D}^\top \mathbf{D})^{-1}\mathbf{P},
\label{eq:hp-closed}
\end{equation}
where $\mathbf{D}$ is the second-difference operator. Larger $\lambda$ implies
a smoother (slower) trend. Hodrick and Prescott (1997) recommend
$\lambda{=}1600$ for quarterly macroeconomic data; on daily financial series
much larger values (e.g.\ $10^4$--$10^5$) are typical.

\paragraph{Causality caveat.} Equation~\eqref{eq:hp-closed} uses the
\emph{entire} series. In a backtest at bar $t$, the historical estimate
$\widehat{\tau}_t^{\text{full}}$ depends on future prices and is therefore
\emph{look-ahead biased}. We address this with a rolling/expanding-window
wrapper (Section~\ref{sec:causal-wrappers}).

\subsection{LOWESS Smoothing}

LOWESS (Cleveland, 1979) fits a low-degree polynomial via locally weighted
least-squares around each target abscissa $x_t$. The weighting uses a
tricube kernel over the nearest \texttt{frac}$\,\cdot\,T$ neighbours. As with
HP, the bare smoother is non-causal in the same global sense: each smoothed
point uses neighbours on both sides. The same rolling-window discipline applies.

\subsection{Exponential Moving Average}

The EMA of span $n$ is recursively
\begin{equation}
\text{EMA}_t \;=\; \alpha\,P_t + (1-\alpha)\,\text{EMA}_{t-1},\qquad \alpha = \frac{2}{n+1}.
\end{equation}
The EMA is \emph{causal by construction}: $\text{EMA}_t$ only depends on
$\{P_s\}_{s\le t}$. It is therefore the natural baseline.

% ============================================================================
\section{Strategies}\label{sec:strategies}
% ============================================================================

Each strategy produces a binary trend state $T_t\in\{-1,+1\}$ on a daily bar.

\begin{table}[h]
\centering
\begin{tabular}{@{}clll@{}}
\toprule
\textbf{\#} & \textbf{Name} & \textbf{Indicator} & \textbf{Tunable parameter(s)} \\
\midrule
S1 & EMA Crossover  & $\operatorname{sign}\!\bigl(\text{EMA}_{n_f} - \text{EMA}_{n_s}\bigr)$ & $(n_f, n_s)$ \\
S2 & EMA Direction  & $\operatorname{sign}\!\bigl(\Delta \text{EMA}_n\bigr)$ & $n$ \\
S3 & HP Direction \emph{(primary)} & $\operatorname{sign}\!\bigl(\Delta \text{HP}_\lambda(P)\bigr)$ & $\lambda$, window $W$ \\
S4 & LOWESS Direction & $\operatorname{sign}\!\bigl(\Delta \text{LOWESS}_{\text{frac}}(P)\bigr)$ & frac, window $W$ \\
\bottomrule
\end{tabular}
\caption{Four trend extractors. All produce a discrete $\pm 1$ trend state.}
\end{table}

\paragraph{Trade rules (common to all variants).}
\begin{itemize}[leftmargin=1.5em,itemsep=2pt]
    \item \textbf{Long entry} when $T_t$ flips from $-1$ to $+1$.
    \item \textbf{Close long and short entry} when $T_t$ flips from $+1$ to $-1$.
    \item \textbf{Always in market}: no flat state in the base strategy.
\end{itemize}

\paragraph{Position sizing --- full capital deployment.}
Sizing uses the course-mandated \texttt{FixedCashSizer} which targets
$\sim\!100{,}000$~TL of notional per individual \texttt{buy()}/\texttt{sell()}
order. A naive implementation that holds a single 100k lot at a time would
leave 90\% of the 1M~TL capital idle, producing misleading single-digit CAGRs.
We therefore fire $\lfloor \text{equity}/100{,}000 \rfloor$ market orders on
each signal flip ($\approx 10$~orders for a fresh 1M account), close any
opposite position first, and hold the stacked lots until the next flip. This
respects the per-trade cap while keeping capital fully invested.

\paragraph{Why HP is designated the primary strategy.}
Per Section~\ref{sec:theory}, the HP-direction rule maps directly onto the
Harris--Yilmaz argument that a low-frequency filter's directional sign is a
clean trend signal. Empirically, LOWESS edges HP at every aggregation level
(Section~\ref{sec:results-h2h}); the difference is small enough that retaining
HP as the principled primary and reporting LOWESS as the head-to-head
benchmark is honest reporting rather than benchmark-swapping.

\subsection{Regime Filter (Phase 2)}\label{sec:regime}

After picking a primary on the index, we deploy it on individual stocks
\emph{gated by the BIST100's own HP-trend}:
\begin{itemize}[leftmargin=1.5em,itemsep=2pt]
    \item \textbf{Long entry} on stock $i$ requires: stock signal long
          \textbf{AND} BIST100 close $>$ HP-filtered BIST100 close.
    \item \textbf{Short entry} on stock $i$ requires: stock signal short
          \textbf{AND} BIST100 close $<$ HP-filtered BIST100 close.
    \item \textbf{Exits} ignore the regime filter (to avoid trapping positions
          when the regime flips intra-trade).
\end{itemize}
The economic story is straightforward top-down filtering: trade with the
index, not against it.

% ============================================================================
\section{Methodology}
% ============================================================================

\subsection{Data and Universe}

\begin{itemize}[leftmargin=1.5em,itemsep=2pt]
    \item \textbf{Phase~1 universe:} BIST100 Index daily close, ticker
          \texttt{XU100.IS} on Yahoo Finance.
    \item \textbf{Phase~2 universe:} 29 BIST100 constituents (cleaned from a
          curated list of 31 well-known blue chips, two failing yfinance ingest
          as ``possibly delisted''). Three additional placeholders
          (\texttt{SODA.IS}, \texttt{KOZAL.IS}, \texttt{KOZAA.IS}) need ticker
          remapping before they can be re-added. Survivorship bias is
          acknowledged but not corrected, per brief.
    \item \textbf{Data fields:} OHLCV in TL, daily frequency.
    \item \textbf{Period:} 2003-01-01 to 2026-05-11 ($\sim$23~years). Panel
          shape: $5997 \times 29$. \texttt{PGSUS} has the shortest history at
          $3273$ bars (later IPO); all 29 clear the
          \texttt{MIN\_HISTORY\_DAYS=750} floor.
\end{itemize}

\paragraph{Currency redenomination and corporate-action fix.}
yfinance does not back-adjust for the 2005-01-03 Turkish TL 1M-to-1
redenomination, leaving a $\sim$1000$\times$ single-day price jump in every
Turkish equity series. The function \texttt{src/data/clean.py::\_adjust\_splits}
detects any single-day price ratio $>3\times$ (or $<1/3\times$) and rescales
all prior OHLC values by that ratio, idempotently. The $3\times$ threshold
(originally $50\times$) also catches the MGROS 2009-08-04 corporate-action gap
(factor $4.04\times$) that previously caused $-14358\%$ MDD blowups on shorts.
The largest legitimate single-day move in the universe is HEKTS $+44\%$, well
below the threshold.

\subsection{Causal Wrappers for HP and LOWESS}\label{sec:causal-wrappers}

The brief's \texttt{hp\_filter(x, lam=1600)} and the
\texttt{statsmodels.lowess(...)} routines both operate on the \emph{whole}
input series at once. Used directly inside a backtest, they leak future data
into past trend values. We therefore wrap both in rolling-window helpers
(\texttt{src/features/hp.py::rolling\_hp\_trend} and
\texttt{src/features/lowess.py::rolling\_lowess\_trend}) that, at each bar
$t$, re-fit on the trailing window $[t-W,t]$ and return only the rightmost
trend value. The bare functions remain available for EDA and plotting; the
backtest pipeline never sees them.

\subsection{Backtest Mechanics}

A single Backtrader Cerebro instance per (strategy, ticker) pair keeps P\&L
attribution clean. The wiring is:

\begin{itemize}[leftmargin=1.5em,itemsep=2pt]
    \item \textbf{Feed} (\texttt{src/backtest/feeds.py}):
          \texttt{PandasData} extended with a precomputed \texttt{trend} line
          and (for the regime variant) a \texttt{regime} line.
    \item \textbf{Strategy} (\texttt{src/backtest/strategies.py}): reads
          \texttt{self.data.trend[0]}, acts on sign flips, sizes orders with
          \texttt{FixedCashSizer(100\_000)}.
    \item \textbf{Runner} (\texttt{src/backtest/runner.py}): orchestrates the
          Cerebro lifecycle, attaches analyzers for per-bar equity and
          per-trade P\&L, returns a structured \texttt{RunResult}.
\end{itemize}

Backtrader's \texttt{next()} with market orders fills the next bar; signals
generated on bar $t$ therefore translate to fills at $t+1$, which is the
correct convention.

\subsection{Performance Metrics}

Reported for every (strategy, variant, ticker) combination:
\begin{multicols}{2}
\begin{enumerate}[leftmargin=1.4em,itemsep=1pt]
    \item Net P/L (TL and \%)
    \item Compound Annual Growth Rate (CAGR)
    \item Annualised volatility
    \item Sharpe ratio (annualised)
    \item Sortino ratio (annualised)
    \item Max drawdown and drawdown duration
    \item Calmar ratio
    \item Win rate, average win, average loss
    \item Number of trades, average holding period
    \item Turnover
\end{enumerate}
\end{multicols}

\subsection{Hyperparameter Optimisation}\label{sec:opt}

\textbf{Phase~1 grid (BIST100 Index)} via
\texttt{src/backtest/sweep.py::default\_specs()}:
\begin{itemize}[leftmargin=1.5em,itemsep=2pt]
    \item \textbf{S1:} $n_f\!\in\!\{5,10,20,30\}$, $n_s\!\in\!\{50,100,200\}$ with $n_f<n_s$.
    \item \textbf{S2:} $n\!\in\!\{10,20,50,100,200\}$.
    \item \textbf{S3:} $\lambda\!\in\!\{100, 1600, 14400, 129600\}$,
          window $W\!\in\!\{252,504,1260\}$.
    \item \textbf{S4:} $\text{frac}\!\in\!\{0.05, 0.1, 0.2, 0.3\}$,
          window same as S3.
\end{itemize}
\textbf{Selection rule:} highest Sharpe, with a minimum-trade floor
($\ge 30$ round trips) to disqualify degenerate configurations.

\textbf{Phase~2 walk-forward:} 3-year training, 1-year test, 1-year step, with
parameter re-selection on each training window. The 5-year ($W=1260$)
configurations are pruned because their warmup leaves too little post-warmup
room for fold-level retraining.

\subsection{Statistical Significance: Monte-Carlo Bootstrap}\label{sec:mc}

For each (strategy, ticker) we generate $N=1000$ random binary signals that
match the empirical:
\begin{enumerate}[leftmargin=1.5em,itemsep=2pt]
    \item run-length distribution (a proxy for trade frequency $\times$
          holding period);
    \item long/short proportion.
\end{enumerate}
Each random signal is then scored on the same prices with the same simulator,
producing a null distribution of Sharpe ratios. The one-sided $p$-value is the
fraction of random Sharpes exceeding the strategy's. The fast vectorised
simulator in \texttt{src/eval/montecarlo.py} is used for both the strategy and
the matched-random pool to keep the comparison apples-to-apples.

\subsection{Portfolio Aggregation}\label{sec:portfolio-method}

The equal-weight cross-sectional portfolio (DESIGN.md milestone M4) is built
by an \texttt{effective\_position} state machine in
\texttt{src/eval/portfolio.py} that mirrors \texttt{TrendStrategy} and
\texttt{TrendRegimeStrategy} exactly, including the regime variant's
\emph{sticky-flat} entry behaviour. Per ticker, this gives a $\{-1,0,+1\}$
held-position series. The portfolio return on bar $t$ is the cross-sectional
mean of per-ticker strategy returns ($\text{pos}_{t-1} \cdot
\Delta\!\ln P_t$), skipping pre-IPO NaNs and counting flat-but-active tickers
as zero contributions.

% ============================================================================
\section{Results}\label{sec:results}
% ============================================================================

\subsection{Phase 1 -- BIST100 Index Sweep}

The 41-configuration grid (run time $\approx 17$~min on a single core, full
deployment) produces the following best-by-Sharpe winners
(\texttt{results/phase1\_sweep.parquet}):

\begin{table}[h]
\centering
\begin{tabular}{@{}llrrrr@{}}
\toprule
\textbf{Strategy} & \textbf{Best params} & \textbf{Sharpe} & \textbf{Trades} & \textbf{CAGR} & \textbf{MDD} \\
\midrule
S4 LOWESS direction          & frac=0.2, $W$=252       & 0.876 & 560 & 19.89\% & $-46.4\%$ \\
S3 HP direction (\emph{primary}) & $\lambda$=14400, $W$=504 & 0.818 & 576 & 17.52\% & $-43.3\%$ \\
S2 EMA direction              & $n$=20                  & 0.742 & 634 & 16.71\% & $-47.1\%$ \\
S1 EMA crossover              & $n_f$=5, $n_s$=100      & 0.719 & 126 & 15.60\% & $-52.8\%$ \\
\midrule
\multicolumn{2}{l}{\emph{Reference: BIST100 buy-and-hold (23.3y)}} & 0.92 & --- & 23.46\% & $-63.4\%$ \\
\bottomrule
\end{tabular}
\caption{Phase-1 winners on the BIST100 Index. The strategies underperform
buy-and-hold on absolute return but cut drawdown by 15--20~percentage points.}
\end{table}

\paragraph{Reading.} Four observations stand out:
\begin{enumerate}[leftmargin=1.5em,itemsep=2pt]
    \item All four strategies are competitive with buy-and-hold in Sharpe
          terms.
    \item LOWESS and HP --- the two low-pass filters --- beat the two EMA
          strategies, consistent with the Harris--Yilmaz hypothesis.
    \item S1 (the classical MA crossover) is the weakest, with the fewest
          trades. Crossover rules wait for a magnitude-difference rather than
          a slope-change, so they enter later and exit later.
    \item Drawdowns are 15--20pp shallower than buy-and-hold but still
          $\sim$45--53\%. This is the cost of always-in-market with shorts.
\end{enumerate}

\subsection{Phase 1 -- Monte-Carlo Significance}

\begin{table}[h]
\centering
\begin{tabular}{@{}lrrrr@{}}
\toprule
\textbf{Strategy} & \textbf{Sharpe} & \textbf{MC mean} & \textbf{MC $q_{95}$} & \textbf{$p$-value} \\
\midrule
S4 LOWESS    & 0.87 & 0.16 & 0.50 & $<0.001$ \\
S2 EMA dir   & 0.86 & 0.19 & 0.54 & $0.002$  \\
S3 HP        & 0.80 & 0.14 & 0.47 & $0.001$  \\
S1 EMA xover & 0.61 & 0.29 & 0.67 & $0.086$  \\
\bottomrule
\end{tabular}
\caption{Run-length-bootstrap Monte-Carlo Sharpe $p$-values, $N{=}1000$.
S2/S3/S4 reject random at $p<0.005$. S1 is only marginal at the 5\% level ---
consistent with the Harris--Yilmaz finding that LF-filter direction
out-performs MA crossover.}
\end{table}

\subsection{Phase 2 -- Per-Stock Full History}

Single-config full-history backtests on 29 constituents, both base and
BIST100-HP regime-gated, run via \texttt{src/eval/run\_phase2.py}:

\begin{table}[h]
\centering
\begin{tabular}{@{}lrrrr@{}}
\toprule
& \textbf{HP base} & \textbf{HP regime} & \textbf{LOWESS base} & \textbf{LOWESS regime} \\
\midrule
mean Sharpe   & 0.419 & 0.372 & 0.454 & 0.421 \\
mean CAGR     & 9.82\% & 7.56\% & 11.37\% & 9.11\% \\
mean MDD      & $-70.4\%$ & $-66.6\%$ & $-72.9\%$ & $-65.7\%$ \\
total trades  & 16{,}717 & 11{,}630 & 17{,}226 & 11{,}592 \\
Sharpe$>$0    & 26/29 & 28/29 & 27/29 & 28/29 \\
\bottomrule
\end{tabular}
\caption{Phase-2 per-stock summary on 29 redenomination-adjusted BIST100
constituents.}
\end{table}

\paragraph{Reading.} The regime filter cuts trade count by $\sim$30\% (entries
gated, exits not), trims mean MDD by 2--5~pp, and lifts the
Sharpe-positive count from 26/29 to 28/29 --- but it does \emph{not} uniformly
raise Sharpe (base edges regime on the mean). It rescues weak tickers (BIMAS,
ENKAI, ARCLK) while hurting some strong ones (TKFEN, TOASO, VESTL, HEKTS).
The headline interpretation: the regime acts as a \emph{variance reducer},
not a \emph{Sharpe enhancer}.

\subsection{Phase 2 -- Per-Stock Monte Carlo}

\begin{table}[h]
\centering
\begin{tabular}{@{}lrr@{}}
\toprule
& \textbf{HP} & \textbf{LOWESS} \\
\midrule
$p<0.01$            & 8/29  & 10/29 \\
$p<0.05$            & 16/29 & 18/29 \\
mean strat Sharpe   & 0.436 & 0.465 \\
\bottomrule
\end{tabular}
\caption{Per-stock MC Sharpe $p$-values, $N{=}1000$, base variant. Regime
variant is intentionally not MC-tested because the gated-entry / sticky-flat
logic is path-dependent and would be mis-stated by the fast simulator.}
\end{table}

Roughly 55\% of constituents reject random at 5\% under HP; 62\% under LOWESS.
At the 1\% level the counts are 28\% and 34\% respectively.

\subsection{Phase 2 -- Equal-Weight Portfolio}\label{sec:portfolio-results}

The cross-sectional equal-weight portfolio (Section~\ref{sec:portfolio-method})
is the headline DESIGN.md~M4 number:

\begin{table}[h]
\centering
\begin{tabular}{@{}lrrrr@{}}
\toprule
\textbf{Metric} & \textbf{HP base} & \textbf{HP regime} & \textbf{LOWESS base} & \textbf{LOWESS regime} \\
\midrule
Sharpe              & 0.895 & 0.823 & 1.049 & 0.958 \\
Sortino             & 1.07  & 0.98  & ---   & ---   \\
CAGR                & 15.29\% & 11.94\% & 18.15\% & 13.74\% \\
Ann.~volatility     & 17.7\% & 15.1\% & 17.3\% & 14.6\% \\
Max drawdown        & $-26.8\%$ & $-25.1\%$ & $-28.5\%$ & $-25.5\%$ \\
Calmar              & 0.57  & 0.48  & --- & --- \\
Total return (24y)  & 29.6$\times$ & 14.7$\times$ & 52.9$\times$ & 21.4$\times$ \\
Mean active tickers & 24/29 & 19/29 & 24/29 & 19/29 \\
\bottomrule
\end{tabular}
\caption{Equal-weight cross-sectional portfolio metrics. The portfolio Sharpe
($\sim$0.9 for HP, $\sim$1.05 for LOWESS) is roughly $2\times$ the mean
per-ticker Sharpe ($\sim$0.42--0.46) --- the diversification benefit of
cross-sectional aggregation.}
\end{table}

\paragraph{Reading.}
\begin{itemize}[leftmargin=1.5em,itemsep=2pt]
    \item LOWESS portfolio Sharpe $1.05$ exceeds BIST100 buy-and-hold Sharpe
          $0.92$. HP portfolio Sharpe $0.90$ is essentially at parity with
          buy-and-hold.
    \item Total nominal return is $\sim$22\% of buy-and-hold for HP and
          $\sim$39\% for LOWESS --- i.e.\ the strategies give up substantial
          nominal upside in exchange for $\sim$2$\times$ shallower drawdowns
          ($-27\%$ vs $-63\%$).
    \item The regime variant lowers volatility (by $\sim$2.6pp for HP,
          $\sim$2.7pp for LOWESS) and tightens MDD marginally, but at the cost
          of $\sim$3--4pp of CAGR and $\sim$0.07--0.09 of Sharpe.
\end{itemize}

\subsection{Phase 2 -- Walk-Forward (S3 HP)}

A 6-configuration grid (\(\lambda\!\in\!\{1600,14400,129600\}\) \(\times\)
\(W\!\in\!\{252,504\}\)), with 3-year train, 1-year test, 1-year step,
produces 477 fold rows per variant:

\begin{table}[h]
\centering
\begin{tabular}{@{}lrrrr@{}}
\toprule
& \textbf{mean test Sharpe} & \textbf{median} & \textbf{\% folds $>0$} & \textbf{mean test MDD} \\
\midrule
Base               & 0.358 & 0.369 & 66.3\% & $-25.4\%$ \\
BIST100-HP regime  & 0.294 & 0.295 & 63.7\% & $-25.6\%$ \\
\bottomrule
\end{tabular}
\caption{Walk-forward summary, S3~HP, 29~tickers.}
\end{table}

\paragraph{Parameter selection.} Across base folds, $W{=}252$ is picked in
$466/476$ folds; $\lambda$ splits $192 / 141 / 133$ across
$\{129600, 14400, 1600\}$. The slow $\lambda$ $\times$ long $W$
combinations are essentially unused --- they should be trimmed in any
next-iteration grid.

\paragraph{Best / worst tickers.} Best by mean fold Sharpe: VESTL~0.86,
ULKER~0.72, TKFEN~0.68, HEKTS~0.67, KRDMD~0.67, KCHOL~0.66, ISCTR~0.63.
Worst: BIMAS~$-0.28$, ENKAI~$-0.23$, TCELL~0.00, SAHOL~0.10, ASELS~0.11.

\subsection{HP vs LOWESS Head-to-Head}\label{sec:results-h2h}

\begin{table}[h]
\centering
\begin{tabular}{@{}lrrr@{}}
\toprule
\textbf{Metric}                  & \textbf{S3 HP} & \textbf{S4 LOWESS} & \textbf{$\Delta$ (L--H)} \\
\midrule
mean Sharpe (base)               & 0.419 & 0.454 & $+0.035$ \\
mean Sharpe (regime)             & 0.372 & 0.421 & $+0.049$ \\
mean MDD (base)                  & $-70.4\%$ & $-72.9\%$ & $-2.5$pp \\
mean MDD (regime)                & $-66.6\%$ & $-65.7\%$ & $+0.9$pp \\
total trades (base)              & 16{,}717 & 17{,}226 & $+509$ \\
Sharpe$>$0 (base)                & 26/29 & 27/29 & $+1$ \\
\midrule
MC $p<0.01$                      & 8/29  & 10/29 & $+2$ \\
MC $p<0.05$                      & 16/29 & 18/29 & $+2$ \\
mean MC Sharpe                   & 0.436 & 0.465 & $+0.029$ \\
\midrule
Portfolio Sharpe (base)          & 0.895 & 1.049 & $+0.154$ \\
Portfolio Sharpe (regime)        & 0.823 & 0.958 & $+0.135$ \\
Portfolio CAGR (base)            & 15.29\% & 18.15\% & $+2.86$pp \\
Portfolio CAGR (regime)          & 11.94\% & 13.74\% & $+1.80$pp \\
Portfolio total return (base)    & 29.6$\times$ & 52.9$\times$ & $+78\%$ \\
\bottomrule
\end{tabular}
\caption{Head-to-head HP vs LOWESS across every aggregation level.}
\end{table}

\paragraph{Verdict.} LOWESS dominates HP at every aggregation level ---
per-stock ($+0.04$ to $+0.06$ Sharpe), MC (more tickers significant at
both 1\% and 5\%), and portfolio ($+0.15$ Sharpe, $+2.8$pp CAGR,
$\sim$80\% more total return). The original choice of S3~HP as the primary
strategy was the theoretically motivated one (Harris--Yilmaz LF-filter
argument); the writeup reports \textbf{HP as primary on principle and LOWESS
as the better empirical alternative} rather than swapping winners. Both reject
random handily; the gap is real but modest in significance terms
($8 \to 10$ tickers at $p<0.01$, $16 \to 18$ at $p<0.05$).

% ============================================================================
\section{Findings and Discussion}
% ============================================================================

\subsection{Summary of Key Findings}

\begin{enumerate}[leftmargin=1.5em,itemsep=4pt]

\item \textbf{Low-frequency filters beat moving-average crossovers}, as
predicted by Harris and Yilmaz. On the BIST100 Index, S3~HP and S4~LOWESS
both produce Sharpe $>0.8$ and reject the matched-random null at $p<0.005$,
while S1~EMA crossover only reaches Sharpe $0.61$ and is marginal at $p=0.086$.

\item \textbf{Trend signal direction (S2/S3/S4) is the consistent winner},
suggesting that the act of \emph{differencing the smoother} --- rather than
\emph{comparing two smoothers} --- is where the edge lives.

\item \textbf{LOWESS narrowly outperforms HP empirically at every aggregation
level} (per-stock, MC, portfolio). The gap is real but modest in statistical
terms.

\item \textbf{The regime filter is a variance reducer, not a Sharpe
enhancer.} Across both HP and LOWESS, gating with the BIST100-HP filter
cuts trade count $\sim$30\% and trims MDD by 2--5pp, but reduces mean Sharpe
slightly. It rescues the worst-of-base tickers (THYAO, BIMAS, ENKAI) while
hurting some of the best (TKFEN, TOASO, VESTL).

\item \textbf{Cross-sectional diversification roughly doubles Sharpe.} The
$\sim$0.42 mean per-ticker Sharpe becomes a $\sim$0.90 portfolio Sharpe ---
the standard diversification dividend from 29 imperfectly correlated bets.

\item \textbf{The strategy is a drawdown reducer, not a return enhancer.}
Versus BIST100 buy-and-hold (CAGR $23.5\%$, MDD $-63.4\%$), the LOWESS
portfolio (CAGR $18.2\%$, MDD $-28.5\%$) gives up $\sim$5pp of CAGR for
roughly half the drawdown. The HP portfolio gives up $\sim$8pp of CAGR for
the same drawdown profile.

\item \textbf{Out-of-sample Sharpe materially lower than in-sample.} The
3y/1y/1y walk-forward yields a $\sim$0.36 mean test Sharpe per ticker, well
below the $0.82$ in-sample index Sharpe. This is the honest cost of fold-level
re-estimation on noisy individual names rather than the smoother index.

\item \textbf{Parameter robustness.} Across walk-forward folds, the short
window ($W{=}252$) wins 466/476 times. $\lambda$ is much less determined:
all three values appear meaningfully (133, 141, 192 selections). The grid
should be trimmed in any next iteration.

\end{enumerate}

\subsection{The Sharpe vs CAGR Trade-Off Frontier}

\begin{table}[h]
\centering
\begin{tabular}{@{}lrrr@{}}
\toprule
\textbf{Approach} & \textbf{Sharpe} & \textbf{CAGR} & \textbf{MDD} \\
\midrule
BIST100 buy-and-hold        & 0.92 & 23.46\% & $-63.4\%$ \\
HP portfolio (base)         & 0.90 & 15.29\% & $-26.8\%$ \\
HP portfolio (regime)       & 0.82 & 11.94\% & $-25.1\%$ \\
LOWESS portfolio (base)     & 1.05 & 18.15\% & $-28.5\%$ \\
LOWESS portfolio (regime)   & 0.96 & 13.74\% & $-25.5\%$ \\
\bottomrule
\end{tabular}
\caption{Risk-adjusted comparison. Only LOWESS-base produces a Sharpe
ratio meaningfully above buy-and-hold; all four portfolio variants cut MDD by
$\sim$35--40pp.}
\end{table}

The natural reading: a leverage-tolerant investor with a moderate $\sim$30\%
drawdown budget can scale the LOWESS-base portfolio by $\sim$2.3$\times$ to
match buy-and-hold's MDD, in which case its expected CAGR would rise to
$\sim$42\% --- well above buy-and-hold's $23\%$. (This is a back-of-envelope
calculation; an honest scaled version would need a re-run with margin
financing costs and short-borrow costs explicitly modelled.)

\subsection{Why the Regime Filter Underperforms on Sharpe}

Asymmetric gating is the proximate culprit: the filter blocks \emph{entries}
but not \emph{exits}. This produces \emph{sticky-flat} stretches where the
signal has flipped (in favour of a new direction), but the regime disagrees,
so the strategy sits in cash --- missing the early part of the new move. When
the regime later confirms, the strategy enters late at a worse price. The
losses on early-entry days that the filter \emph{would have caught} are
typically smaller than the missed early gains on the days the filter blocks.

\subsection{Comparison with Harris--Yilmaz (2008)}

Harris and Yilmaz study daily FX, not equities. Two systematic differences
matter:

\begin{itemize}[leftmargin=1.5em,itemsep=2pt]
    \item \textbf{Drift sign.} Equities have a positive unconditional drift,
          FX pairs do not. Always-in-market shorts on equities therefore bleed
          against the drift even on good signals. This explains the large
          per-stock MDDs and the strategy's underperformance versus
          buy-and-hold on absolute return.
    \item \textbf{Volatility regimes.} Turkish equities exhibit the FX-like
          regime-of-regimes structure (currency crises, 2018 lira crisis, etc.)
          but with additional equity-specific factors (sector cycles, corporate
          actions). The HP/LOWESS direction signal captures the
          regime-of-regimes well; sector-specific factors are uncaptured.
\end{itemize}

The structural Harris--Yilmaz conclusion --- LF-filter direction beats MA
crossover --- replicates in our setting. The magnitude of the edge is smaller,
which is what one would expect when transferring an FX result to an equity
universe with positive drift and shortable but costly individual names.

% ============================================================================
\section{Improvement Ideas and Future Work}
% ============================================================================

The following are concrete, prioritised improvement directions. Each is scoped
small enough to be tackled as an independent piece of work; none blocks the
others.

\subsection{Data and Universe}

\paragraph{IM1 -- Full official BIST100 membership.} The current 29-name
universe is a curated subset of well-known blue chips. Swapping in the
official current 100-name membership (subtracting whichever names still fail
yfinance) would treble the cross-sectional sample, improve diversification,
and let us report per-sector breakdowns honestly. \emph{Estimated effort:}
1--2~days, mostly ticker remapping for the names that still fail ingest.

\paragraph{IM2 -- Point-in-time BIST100 membership (no survivorship).} The
brief permits survivorship bias, but the most defensible upgrade is to use
\emph{historical point-in-time} membership: at each rebalance date, only
trade names that were in the index at that time. Bloomberg and Refinitiv ship
this; constructing it from public sources is laborious but not infeasible.
\emph{Estimated effort:} 1--2~weeks if done from scratch, 1~day if a vendor
file is available.

\paragraph{IM3 -- Adjusted-close vendor data.} The
\texttt{\_adjust\_splits($3\times$)} threshold catches the 2005 redenomination
and the MGROS 2009 gap; but a single threshold is fragile. Using a
proper corporate-actions feed (vendor adjusted-close with audited split,
dividend and rights-issue factors) would eliminate the threshold-tuning
fragility entirely.

\paragraph{IM4 -- Currency-translated returns.} All returns are in TL. The TL
has depreciated $\sim$90\% in USD terms over the sample. A USD-denominated
view would expose to a foreign investor what fraction of the local
Sharpe survives FX translation. This is a $\sim$1-day add-on (multiply by
USDTRY, compute Sharpe).

\subsection{Signal Engineering}

\paragraph{IM5 -- HP + LOWESS signal ensemble.} The two LF-filter signals are
correlated but not identical (Phase-1 Sharpe gap of $0.06$, per-stock Sharpe
gap of $0.04$). A simple AND-gate (trade only when both agree) or weighted
vote could be tested cheaply. \emph{Hypothesis:} ensembling cuts whipsaws at
the cost of a slightly slower entry, with net Sharpe improvement.

\paragraph{IM6 -- Adaptive $\lambda$.} The walk-forward shows
$\lambda$-selection is unstable across folds (no single $\lambda$ dominates,
unlike $W{=}252$). A volatility-adaptive $\lambda(t) = \lambda_0 / \sigma_t^2$
(more smoothing in high-vol regimes) is a one-line change and would test
whether the instability is genuine or an artefact of grid coarseness.

\paragraph{IM7 -- Slope magnitude, not just sign.} The current rules use only
$\operatorname{sign}(\Delta \tau_t)$. Using $\Delta \tau_t / \sigma_t$ as a
\emph{continuous} signal that scales position size would give a smoother
allocation and might reduce whipsaw losses. This is a meaningful re-design ---
the brief's binary-trend convention would need to be relaxed.

\paragraph{IM8 -- HP residual as a contrarian overlay.} The high-frequency
residual $c_t = P_t - \tau_t$ is by construction mean-reverting. A short-term
mean-reversion overlay using $c_t$ could be added as an orthogonal alpha; the
two should diversify well.

\paragraph{IM9 -- Kalman / state-space trend.} The HP filter is a special case
of a local-linear-trend Kalman model with specific signal-to-noise ratio.
Estimating the SNR from data (rather than fixing $\lambda$) is a principled
upgrade with virtually no extra implementation cost on top of an
\texttt{statsmodels} \texttt{UnobservedComponents} call.

\subsection{Position Sizing and Risk Management}

\paragraph{IM10 -- Volatility targeting.} Current sizing is fixed-cash
(100k~TL per order). A natural improvement is to size to a constant
\emph{volatility} target: $\text{notional}_t \propto 1/\hat{\sigma}_t$. This
is standard in CTA practice; it should reduce realised vol drift and improve
risk-adjusted return at the portfolio level. The brief's
\texttt{FixedCashSizer} could be replaced with a volatility-targeted sizer
without otherwise touching the pipeline.

\paragraph{IM11 -- Cross-sectional risk-parity weighting.} The portfolio is
currently equal-weight. Risk-parity (weight $\propto 1/\sigma_i$) would
reduce the dominance of high-vol names in the realised portfolio variance.
This is a one-line change in
\texttt{src/eval/run\_phase2\_portfolio.py}.

\paragraph{IM12 -- Drawdown stop-out.} Per-ticker MDDs of $-50\%$ to $-90\%$
are the cost of always-in-market shorts on names with positive drift. A
per-ticker drawdown stop (e.g.\ exit and wait for the next signal flip if
the rolling-1y drawdown exceeds 40\%) would directly attack the worst
realised paths. Requires careful handling so it doesn't degenerate into
trend-chasing.

\paragraph{IM13 -- Position-level concentration limits.} Single-name
concentration is currently bounded only by capital. An explicit limit
($\le 10\%$ of equity per name) would protect against catastrophic single-name
events.

\subsection{Backtest Realism}

\paragraph{IM14 -- Transaction costs.} The brief mandates zero commissions,
which the headline backtest respects. A sensitivity analysis with realistic
BIST commissions (typically 5--15~bps round-trip), bid--ask spread (5--20~bps
depending on liquidity), and short-borrow cost (highly name-dependent;
typically 1--5\% annualised for liquid names, much higher for the small-cap
short side) would tell us how much of the Sharpe survives. Expected impact:
$\sim$0.1--0.3 Sharpe loss for the high-trade-count strategies.

\paragraph{IM15 -- Slippage model.} Market orders fill at the next-bar open
in Backtrader's default \texttt{Cheat-On-Open=False} mode. A more honest
fill model (e.g.\ a fraction of average daily volume capped at 5\% of ADV)
would surface any liquidity-constrained names. PGSUS and ARCLK are
candidates.

\paragraph{IM16 -- Borrow availability.} Shorts on Turkish blue chips are
generally feasible but not always cheap. A two-tier model (block the short
if borrow $>$ X bps, scale down if borrow $\in$ [Y,X]) would be the right
realism upgrade. Out of scope of the brief but should be acknowledged.

\subsection{Statistical Testing}

\paragraph{IM17 -- Regime-variant per-stock MC.} The current fast Monte-Carlo
simulator cannot model the regime variant's sticky-flat behaviour and so the
per-stock MC table is base-only. Building
\texttt{monte\_carlo\_sharpe\_regime(prices, signal, regime, n\_iter)} on top
of the existing \texttt{effective\_position} state machine in
\texttt{src/eval/portfolio.py} is a $\sim$50-line change and would close the
methodological gap noted in DESIGN.md~\S10.4.

\paragraph{IM18 -- Portfolio-level Monte Carlo.} The portfolio Sharpes
($0.90$ HP / $1.05$ LOWESS) deserve their own significance test against a
matched-random \emph{portfolio}. Per MC iteration: draw an independent
matched-random signal per ticker, build the effective-position panel,
equal-weight-average, compute portfolio Sharpe. With $N=500$ the resolution
at $p<0.01$ is fine and the runtime is tractable.

\paragraph{IM19 -- Stationary block bootstrap.} The current run-length match
preserves trade-frequency moments but assumes the run-length distribution
itself is stationary. Politis and Romano's stationary block bootstrap is the
standard upgrade, blocking on the price series itself rather than the signal.

\paragraph{IM20 -- White's reality check / SPA test.} With four strategies and
multiple parameter configurations evaluated on the same series, the
single-strategy Sharpe $p$-values are subject to multiple-testing inflation.
White's (2000) reality check and Hansen's SPA test (2005) are the standard
corrections; both have public Python implementations.

\subsection{Walk-Forward and Robustness}

\paragraph{IM21 -- LOWESS walk-forward.} The walk-forward driver currently
hardcodes \texttt{hp\_direction\_signal}. Adding a \texttt{--strategy lowess}
arm and running it ($\sim$25--40 minutes on the full universe) would close
the HP-vs-LOWESS head-to-head out-of-sample. Recommended for the next
iteration; not strictly required by the brief.

\paragraph{IM22 -- Trim the WF grid.} $W{=}504$ is essentially never picked
(10/476 folds); dropping it halves the WF runtime with no measurable loss.

\paragraph{IM23 -- Combinatorial purged cross-validation (Lopez de Prado).}
Walk-forward is a single time-ordered split, vulnerable to a single bad
fold. Purged $k$-fold CV with embargo is the modern standard for financial
time-series CV. Implementation effort: $\sim$1--2 days.

\paragraph{IM24 -- Deflated Sharpe ratio.} Lopez de Prado's \emph{deflated
Sharpe ratio} (2014) corrects for both selection bias (we tried 41 Phase-1
configurations) and the higher moments of returns. Should be reported
alongside the raw Sharpe in any updated table.

\subsection{Engineering and Reproducibility}

\paragraph{IM25 -- Feature parquet cache.} DESIGN.md~\S5.2 planned a
parameter-keyed \texttt{data/features/} cache for indicator parquets and
explicitly deferred it. With the LOWESS portfolio taking $\sim$25 minutes
(vs $\sim$1 minute for HP) because each rolling refit cannot be
pre-factored, the cache becomes worthwhile as soon as we extend the LOWESS
grid materially.

\paragraph{IM26 -- Restore the venv.} The repository's \texttt{.venv} was
created under the old path \texttt{\textasciitilde/ec532-project} (the
repo was renamed). Console scripts with hard-coded shebangs (\texttt{pip},
\texttt{jupyter}, \texttt{nbconvert}) fail with \emph{bad interpreter}.
Workaround in use: invoke via \texttt{.venv/bin/python -m <module>}.
A clean rebuild of the venv would resolve this permanently.

\paragraph{IM27 -- Git initialisation.} The project is not yet a git
repository. A \texttt{git init} + initial commit would enable per-run
result reproducibility hashes in the parquet metadata.

\paragraph{IM28 -- Test coverage for causality.} The most consequential
methodological pitfall in this project is HP/LOWESS look-ahead bias. A unit
test that compares \texttt{rolling\_hp\_trend(P, lam, W)[t]} against
\texttt{hp\_filter(P[:t+1], lam)[-1]} would catch any future regression
where the wrapper accidentally peeks ahead.

\subsection{Reporting Extensions}

\paragraph{IM29 -- Per-sector slice.} \texttt{config.py} has placeholder
sector comments. Adding a sector tag column and a \texttt{groupby} in the
portfolio aggregator would give per-sector Sharpe tables (banks, holdings,
transport, energy, industrials, etc.). A 1-hour change.

\paragraph{IM30 -- Regime-conditional return tables.} Decomposing portfolio
returns into ``BIST100 trending up'' vs ``BIST100 trending down'' weeks
would directly visualise where the strategy makes its money. Useful in the
deck.

\paragraph{IM31 -- Drawdown attribution.} For each MDD episode $>10\%$ in
the portfolio curve, attribute the loss to (a) underlying market move,
(b) signal-side losses, (c) regime gating. This is the diagnostic-style
table that gives a reader confidence in the variance-reducer interpretation.

% ============================================================================
\section{Limitations}
% ============================================================================

\begin{enumerate}[leftmargin=1.5em,itemsep=4pt]
    \item \textbf{Survivorship bias.} The 29-name universe is the
          \emph{current} membership snapshot. Names that entered then left
          the index, and names that delisted in the sample, are absent. The
          brief permits this but a point-in-time membership reconstruction
          (IM2) is the right next step.

    \item \textbf{Zero transaction costs.} Mandated by the brief, but the
          strategies trade 16{,}000--17{,}000 round-trips on the
          29-ticker universe over 24 years. At a realistic
          10~bps round-trip cost, the unhedged Sharpe loss would be
          $\sim$0.1--0.2 (IM14).

    \item \textbf{No short-borrow cost.} Equally mandated by the brief;
          arguably the more important omission given that the strategy spends
          $\sim$50\% of its life short.

    \item \textbf{No corporate-action vendor feed.} Splits, rights issues
          and dividends are not properly back-adjusted; we use a single
          $>3\times$ price-jump heuristic in \texttt{\_adjust\_splits}.
          Robust to the 2005 redenomination and the MGROS 2009 gap, but
          could miss smaller corporate actions (IM3).

    \item \textbf{Single primary strategy.} S3~HP is the primary on
          theoretical grounds; LOWESS is the empirical winner. A
          model-selection upgrade (cross-validated Sharpe on a holdout fold)
          would resolve the choice on data rather than priors.

    \item \textbf{Backtrader \texttt{multiprocessing} not used.} The brief
          warns this optimiser can be flaky; we used sequential loops. With
          a feature cache (IM25) this becomes a non-issue.

    \item \textbf{Single-window walk-forward.} The 3y/1y/1y split is one
          specific choice; combinatorial purged CV (IM23) would test
          robustness to that choice.

    \item \textbf{Sharpe-only selection criterion.} Hyperparameter selection
          uses only Sharpe (with a min-trade floor). Selection by Calmar,
          Sortino, or a multi-objective rule would likely pick different
          parameters.
\end{enumerate}

% ============================================================================
\section{Conclusion}
% ============================================================================

We implemented a four-variant trend-following strategy on Turkish equities
inspired by Harris and Yilmaz (2008), with full respect for the EC581 course
brief's constraints: Backtrader engine, 1M~TL initial capital, 100k~TL
per-trade \texttt{FixedCashSizer}, zero commissions, daily bars, and a
matched-random Monte-Carlo $p$-value with $N{=}1000$.

The headline empirical findings are:

\begin{itemize}[leftmargin=1.5em,itemsep=2pt]
    \item \textbf{The Harris--Yilmaz hypothesis replicates.} Low-frequency
          filter direction (S3~HP, S4~LOWESS) beats classical MA crossover
          (S1) on the BIST100 Index, both in raw Sharpe and in Monte-Carlo
          significance.
    \item \textbf{HP is the principled primary; LOWESS is the empirical
          winner.} The gap is small (0.04--0.15 Sharpe across aggregation
          levels) but consistent; we report both.
    \item \textbf{The cross-sectional portfolio doubles per-ticker Sharpe.}
          From $\sim$0.42 mean per-stock to $\sim$0.90 (HP) and $\sim$1.05
          (LOWESS) portfolio --- a textbook diversification dividend.
    \item \textbf{The regime filter is a variance reducer.} It cuts trade
          count $\sim$30\%, tightens MDD slightly, but slightly reduces mean
          Sharpe. It rescues weak tickers and clips strong ones.
    \item \textbf{The strategy gives up nominal return for drawdown
          reduction.} HP/LOWESS portfolios produce $\sim$22--39\% of
          buy-and-hold's total return at $\sim$42\% of buy-and-hold's MDD.
\end{itemize}

The most impactful improvement directions are: (i)~adding transaction-cost
and short-borrow realism (IM14--16), (ii)~upgrading the universe to
point-in-time BIST100 membership (IM2), (iii)~implementing the regime-variant
and portfolio-level Monte-Carlo (IM17--18), and (iv)~moving to
volatility-targeted sizing and risk-parity weighting (IM10--11). Each is a
1-day to 1-week piece of work and individually addresses a distinct
limitation in the current methodology.

% ============================================================================
\section*{References}
\addcontentsline{toc}{section}{References}
% ============================================================================

\begin{description}[leftmargin=2em,labelindent=0pt,style=nextline]
    \item[Asness, Moskowitz, Pedersen (2013).] Value and Momentum Everywhere.
          \emph{Journal of Finance} 68(3), 929--985.
    \item[Cleveland, W.S.\ (1979).] Robust Locally Weighted Regression and
          Smoothing Scatterplots. \emph{Journal of the American Statistical
          Association} 74(368), 829--836.
    \item[Hansen, P.R.\ (2005).] A Test for Superior Predictive Ability.
          \emph{Journal of Business and Economic Statistics} 23(4), 365--380.
    \item[Harris, R.D.F.\ and Yilmaz, F.\ (2008/2009).] A Momentum Trading
          Strategy Based on the Low Frequency Component of the Exchange Rate.
          \emph{Journal of Banking and Finance}.
    \item[Hodrick, R.J.\ and Prescott, E.C.\ (1997).] Postwar U.S.\ Business
          Cycles: An Empirical Investigation. \emph{Journal of Money, Credit
          and Banking} 29(1), 1--16.
    \item[Lopez de Prado, M.\ (2014).] The Deflated Sharpe Ratio: Correcting
          for Selection Bias, Backtest Overfitting and Non-Normality.
          \emph{Journal of Portfolio Management} 40(5), 94--107.
    \item[Lopez de Prado, M.\ (2018).] \emph{Advances in Financial Machine
          Learning}. Wiley.
    \item[Moskowitz, Ooi, Pedersen (2012).] Time Series Momentum.
          \emph{Journal of Financial Economics} 104(2), 228--250.
    \item[Politis, D.\ and Romano, J.\ (1994).] The Stationary Bootstrap.
          \emph{Journal of the American Statistical Association} 89(428),
          1303--1313.
    \item[White, H.\ (2000).] A Reality Check for Data Snooping.
          \emph{Econometrica} 68(5), 1097--1126.
    \item[Y\"uksel, A.\ (2026).] EC581 Project Brief --- Algorithmic Trading
          and Quantitative Strategies. Bogazici University.
\end{description}

% ============================================================================
\appendix
\section{Code Structure}\label{app:code}
% ============================================================================

\begin{lstlisting}[basicstyle=\ttfamily\footnotesize,language=,frame=single]
ec581-project/
|-- DESIGN.md                  # design document (full spec)
|-- CLAUDE.md                  # development memo
|-- NB_Projects.pdf            # course brief
|-- config.py                  # paths, seeds, capital, universe
|-- requirements.txt
|-- data/                      # gitignored
|   |-- raw/                   # yfinance ingest, immutable
|   |-- clean/                 # cleaned, redenomination-adjusted
|   `-- panel/                 # wide prices/volumes + bist100 OHLCV
|-- results/                   # parquet outputs (sweep, MC, portfolio, WF)
|-- notebooks/
|   |-- 01_eda.ipynb
|   |-- 02_index_strategy_comparison.ipynb
|   |-- 03_walk_forward.ipynb
|   `-- 04_regime_filter.ipynb
`-- src/
    |-- data/                  # ingest, clean, panel, build
    |-- features/              # ema, hp, lowess, regime (causal wrappers)
    |-- backtest/              # sizers, feeds, strategies, runner, sweep
    `-- eval/                  # metrics, montecarlo, walkforward,
                               # portfolio, strategies registry,
                               # CLIs: run_phase2[_mc|_portfolio],
                               #       run_walkforward
\end{lstlisting}

% ============================================================================
\section{Reproducibility}\label{app:repro}
% ============================================================================

\begin{enumerate}[leftmargin=1.5em,itemsep=2pt]
    \item Build data: \texttt{.venv/bin/python -m src.data.build}
          (full universe) or
          \texttt{.venv/bin/python -m src.data.build --smoke}
          (M1 subset).
    \item Phase-1 sweep:
          \texttt{.venv/bin/python -m src.backtest.sweep}
          $\to$ \texttt{results/phase1\_sweep.parquet}.
    \item Phase-1 Monte Carlo:
          \texttt{.venv/bin/python -m src.eval.montecarlo}
          $\to$ \texttt{results/phase1\_mc.parquet}.
    \item Phase-2 per stock:
          \texttt{.venv/bin/python -m src.eval.run\_phase2
          --strategy \{hp,lowess\}}.
    \item Phase-2 per-stock MC:
          \texttt{.venv/bin/python -m src.eval.run\_phase2\_mc
          --strategy \{hp,lowess\} --n-iter 1000}.
    \item Phase-2 portfolio:
          \texttt{.venv/bin/python -m src.eval.run\_phase2\_portfolio
          --strategy \{hp,lowess\}}.
    \item Phase-2 walk-forward:
          \texttt{.venv/bin/python -m src.eval.run\_walkforward}
          (HP only as of this writing).
    \item Notebook re-execution:
          \texttt{.venv/bin/python -m jupyter nbconvert --to notebook
          --execute notebooks/<nb>.ipynb --output <nb>.ipynb}.
\end{enumerate}

All MC and bootstrap seeds are sourced from \texttt{config.RANDOM\_SEED=42}.

% ============================================================================
\end{document}
% ============================================================================
